\begin{multicols}{2}
    \begin{exercice}{}{}
        Sans calcul, déterminer le signe des produits suivants.
        \begin{enumeratebf}
            \item $A=(-154)\times 97$
            \item $B=(-4)\times (-27)\times 103$
            \item $C=42\times (-2)\times (-87)\times (-9)$
            \item $D=(-74)\times 562\times (-2)\times 5\times (-87)$
            \item $E= (-1)^{2026} = (-1)\times \dots \times (-1)$ (2026 facteurs)
        \end{enumeratebf}
        \textit{Il suffit de compter le nombre de facteurs négatifs dans le produit. Si ce nombre est pair, le produit est positif, sinon il est négatif. Par exemple, $F = (-3) \times (-1) \times (-2)$ a 3 facteurs négatifs, donc $F$ est négatif.}
    \end{exercice}

    \begin{exercice}{}{}
        Effectuer les calculs suivants :
        \begin{enumeratebf}
            \item $A = 2,5 \times 2 \times (-1)$
            \item $B = 10 \times (-2) \times (-0,1)$
            \item $C = (-1) \times (-1) \times (-1) \times (-1)$
            \item $D = 4 \times (-2) \times 3 \times 10$
        \end{enumeratebf}
        \textit{On pourra essayer de regrouper certains facteurs pour les calculer à part, puis remplacer l'expression par le résultat. Par exemple, $12 \times 25 \times 4 = 12 \times (25 \times 4) = 12\times 100 = 1200  $}
    \end{exercice}

    \begin{exercice}{}{}
        Effectuer les calculs suivants :
        \begin{enumeratebf}
            \item $A = 3 \times (-4)$
            \item $B = -5,2 \times 6$
            \item $C = -7 \times (-2)$
            \item $D = -2,5 \times (-3) + 4$
        \end{enumeratebf}
    \end{exercice}

    \begin{exercice}{}{}
        Effectuer les calculs suivants, en posant les multiplications :
        \begin{enumeratebf}
            \item $A = 23 \times 5$
            \item $B = 12,5 \times 4$
            \item $C = 7,2 \times 3$
            \item $D = 15 \times 6$ 
        \end{enumeratebf}
    \end{exercice}
\end{multicols}